\documentclass[10pt]{article}
\usepackage[margin=0.72in]{geometry}
\usepackage{graphicx}
\usepackage[hidelinks]{hyperref}
\usepackage{xcolor}
\setlength{\parindent}{0pt}
\setlength{\parskip}{0.45em}
\pagestyle{empty}

\newcommand{\pending}[1]{\textcolor{red!70!black}{\textbf{[#1]}}}

\begin{document}

\begin{center}
{\Large\textbf{ECE 2390 Milestone: Prohibitory Road-Sign Detection}}\\[-0.1em]
William Sandoval Casas and John Ferguson \hfill November 13, 2026
\end{center}

\textbf{Objective.}
We are establishing a reproducible comparison between a simple classical
color-and-shape detector and the HOG+LDA baseline reported by Houben et al. for
the German Traffic Sign Detection Benchmark (GTSDB). The paper reports 70.33\%
area under its precision--recall curve for prohibitory signs. We have not yet
implemented HOG+LDA and do not claim to reproduce that value.

\textbf{Dataset and split.}
GTSDB contains 900 road images, originally divided into 600 training and 300
evaluation images. We group corrected class IDs
\texttt{0,1,2,3,4,5,7,8,9,10,15,16} as one prohibitory class. The locally
available test ZIP contains 300 images but no ground truth, so it cannot yet be
scored. This milestone will report the deterministic final 20\% of the
annotated training images in filename order (120 images when the full training
package is present) as a \emph{development} result.

\textbf{Implemented pipeline.}
The detector converts each BGR image to HSV, thresholds low- and high-hue red
ranges, applies a $3\times3$ morphological opening and closing, extracts
external contours, and filters candidate boxes by size and aspect ratio. It
operates without access to ground truth. Each candidate is ranked by the
fraction of red-mask pixels inside its box; this is a heuristic score, not a
calibrated probability.

\textbf{Evaluation and results.}
Detections are ranked by confidence and matched one-to-one to ground truth in
the same image. We report conventional AP at IoU 0.50 and 0.60 using all-points
interpolation of the precision envelope; this is not yet the official GTSDB
scoring procedure. Exact split: \pending{insert filenames and image/GT counts
from metrics.json}. AP@0.50: \pending{run annotated development set}.
AP@0.60: \pending{run annotated development set}. At confidence 0.10,
precision/recall: \pending{insert measured operating points}.

\textbf{Observed failure modes.}
\pending{After the real run, summarize examples such as red background clutter,
faded or shadowed signs, small signs removed by morphology, and merged red
regions. Do not infer failures from synthetic tests.}

\textbf{Next steps.}
We will acquire verified annotations, run and inspect the development baseline,
then implement the official GTSDB scoring protocol. After the milestone we will
reproduce HOG+LDA and compare both methods on the same split and protocol.

\vfill
\begin{center}
\IfFileExists{../outputs/milestone_figure.png}{%
  \includegraphics[width=\textwidth]{../outputs/milestone_figure.png}%
}{%
  \fbox{\parbox[c][1.55in][c]{0.94\textwidth}{\centering
  \pending{Run scripts/run\_baseline.py on annotated GTSDB data to generate the
  two-panel milestone figure.}}}%
}
\end{center}

\vspace{-0.6em}
\footnotesize
S. Houben et al., ``Detection of Traffic Signs in Real-World Images: The German
Traffic Sign Detection Benchmark,'' IJCNN, 2013,
\href{https://doi.org/10.1109/IJCNN.2013.6706807}{doi:10.1109/IJCNN.2013.6706807}.

\end{document}

